\documentclass[11pt]{article}
\usepackage[a4paper,margin=23mm]{geometry}
\usepackage{amsmath,amssymb,amsthm,mathtools,bm,microtype}
\usepackage[hidelinks]{hyperref}
\emergencystretch=2em
\newtheorem{theorem}{Theorem}
\newtheorem{proposition}[theorem]{Proposition}
\newtheorem{lemma}[theorem]{Lemma}
\newtheorem{corollary}[theorem]{Corollary}
\newtheorem{definition}[theorem]{Definition}
\theoremstyle{remark}\newtheorem{remark}[theorem]{Remark}
\newcommand{\R}{\mathbb R}\newcommand{\C}{\mathbb C}
\newcommand{\id}{\mathbf 1}\newcommand{\Tr}{\operatorname{Tr}}
\newcommand{\tr}{\operatorname{tr}}\newcommand{\Q}{\mathcal Q}
\newcommand{\W}{\mathcal W}\newcommand{\sxdeg}{\operatorname{sxdeg}}
\newcommand{\Span}{\operatorname{span}}\newcommand{\Newt}{\operatorname{Newt}}
\title{Complete proofs\\Infinite NPA depth at minimal semidefinite complexity}
\author{Anubhav Chaturvedi}\date{28 October 2026}
\begin{document}
\maketitle

Alice and Bob each choose one of two measurements and obtain an outcome
in $\{\pm1\}$. Write $A_0,A_1$ and $B_0,B_1$ for their observables,
and $\Q_{222}$ for the set of quantum correlations in this scenario.
Let $M=\langle A_0\rangle+\langle B_0\rangle$ and
$S=\langle A_0B_0+A_0B_1+A_1B_0-A_1B_1\rangle$ in the two-input,
two-output Bell scenario, and let $\W=\{(M,S):p\in\Q_{222}\}$.
A semidefinite lift describes membership by introducing auxiliary variables
and imposing affine equations and positive semidefinite matrix constraints.
Its block orders are the sizes of these matrices. The semidefinite extension
degree is the smallest possible largest block order, allowing any finite
number of blocks. We give an explicit lift of $\W$ with blocks of orders
five and four and prove that its extension degree is two. Every fixed level
of the Navascu\'es--Pironio--Ac\'in (NPA) hierarchy nevertheless projects
to a strict superset of $\W$.

For each supporting inequality, the quantum maximum can be attained when
Alice's two observables and Bob's two observables have the same relative
angle. Part I proves this reduction and constructs the membership matrices
from pure two-qubit states. It also proves that two such strategies suffice,
treats arbitrary signs and unequal marginal coefficients, and gives
thirty-two equivalent projections under measurement relabellings.
Part II proves that no fixed NPA level is exact. If its error remained
bounded relative to the quantum violation near the local boundary,
the certificate would give an impossible polynomial sum of squares.

\part*{I. The finite representation}
\renewcommand{\thesection}{F\arabic{section}}
\renewcommand{\thesubsection}{\thesection.\arabic{subsection}}
\renewcommand{\theequation}{F\arabic{equation}}
\renewcommand{\thetheorem}{F\arabic{theorem}}
\renewcommand{\theHequation}{F.\arabic{equation}}
\renewcommand{\theHsection}{F.\arabic{section}}
\section{Why equal local measurement angles suffice}

The supporting inequalities of $\W$ have the form $aM+hS\le\lambda$.
Write $\mathrm{CHSH}=A_0B_0+A_0B_1+A_1B_0-A_1B_1$ for the CHSH
operator. We determine when its tilted form is bounded above by $\lambda I$
for every choice of state and measurements.
For \(H=h\,\mathrm{CHSH}+a(A_0+B_0)\), with \(h,a\in\mathbb R\), write
\[
 {\cal C}=\{(\lambda,a,h):\lambda I-H\succeq0
 \text{ for every quantum realization}\}.
\]
Projective dilation and Jordan's lemma reduce this condition to real
two-qubit blocks. With $X,Y,Z$ denoting the Pauli matrices, put
\[
 A_0=B_0=Z,\quad A_1=cZ+rX,\quad B_1=dZ+sX,\qquad
 r^2=1-c^2,\ s^2=1-d^2 .
\]
Here $c,d\in[-1,1]$ are the cosines of the two local measurement angles.
We may take \(r,s\ge0\) by local \(Z\) conjugations. Endpoints include the
one-dimensional blocks. Complex state vectors cause no further condition,
because the Bell matrix is real symmetric.

Set \(H_*=|h|\), \(b=|a|\) and \(L=2H_*+2b\).
Deterministic assignments attain \(L\): choose \(A_0=B_0=\operatorname{sign}(a)\),
and take \(A_1=B_1\) equal to that sign for \(h\ge0\), or its negative for \(h<0\).
At \(a=0\), either sign works. Therefore every admissible \(\lambda\) satisfies
\(\lambda\ge L\).

\paragraph{The determinant controls positivity.}
With \(\Omega=1+c+d-cd\in[-2,2]\), the coefficients in
\[
 \det(uI+\lambda I-H)=u^4+e_1u^3+e_2u^2+e_3u+e_4
\]
are
\begin{align*}
 e_1&=4\lambda,&
 e_2&=6\lambda^2-4a^2-8h^2,\\
 e_3&=4\lambda^3-8\lambda a^2-16\lambda h^2-8ha^2\Omega,
 &e_4&=\det(\lambda I-H).
\end{align*}
The first three are nonnegative for \(\lambda\ge L\).
For \(e_3\), bound its last term below by \(-16H_*b^2\).
The resulting cubic is increasing for \(\lambda\ge L\), and at \(L\) equals
\(16b(b+2H_*)^2\).
Consequently
\[
 \lambda I-H\succeq0\quad\Longleftrightarrow\quad e_4\ge0 .
\]
Indeed, a negative eigenvalue would produce a positive root of a polynomial
whose coefficients are all nonnegative. This argument includes zero determinants
and multiple zero eigenvalues.

\paragraph{At fixed product, minimize the sum.}
Put \(x=(1-c)/2\), \(y=(1-d)/2\), so $(x,y)\in[0,1]^2$, and denote
the determinant $\det(\lambda I-H)$ by \(D(x,y)\).
Direct expansion gives
\begin{align*}
D(x,y)={}&D_0+64h^2(4h^2-a^2)xy(x+y)\\
 &+(32\lambda ha^2+128h^2a^2-256h^4)xy-256h^4x^2y^2,\\
D_0={}&[(\lambda-2h)^2-4a^2](\lambda+2h)^2.
\end{align*}
Only \(x+y\) varies when \(xy\) is fixed. Arithmetic--geometric mean gives
\(x+y\ge2\sqrt{xy}\), and therefore
\begin{equation}
 D(x,y)-D(\sqrt{xy},\sqrt{xy})
 =64h^2(4h^2-a^2)xy(\sqrt x-\sqrt y)^2 .
 \label{eq:diagonal}
\end{equation}
For \(|a|\le |h|\), the right side is nonnegative. Replacing $(x,y)$ by
$(\sqrt{xy},\sqrt{xy})$ stays in $[0,1]^2$ and makes $c=d$.
Thus nonnegative determinants for all equal-angle choices imply nonnegative
determinants for every pair \(c,d\), including the endpoints. For
$\lambda\ge L$, this is equivalent to positivity of the corresponding slacks.

\paragraph{The remaining directions are classical.}
The range \(b\ge H_*\) has exact support value \(L\).
To prove this, set \(C_B=(B_0+B_1)/2\), \(D_B=(B_0-B_1)/2\). Then
\begin{equation}
 4I-(\mathrm{CHSH}+2A_0)
 =(I-C_B-A_1D_B)^2+
 \frac14\sum_{y=0}^1[(I-A_0)(I+B_y)]^2\succeq0 .
 \label{eq:critical}
\end{equation}
Here \(C_B^2+D_B^2=I\) and \(C_BD_B+D_BC_B=0\).
Party exchange and output flips give
\(\sigma\mathrm{CHSH}+2\tau A_0\le4I\) and
\(\sigma\mathrm{CHSH}+2\tau B_0\le4I\), for either sign \(\sigma,\tau\).
Average these inequalities, multiply by \(H_*\), and add
\((b-H_*)[2I-\tau(A_0+B_0)]\succeq0\).
The resulting bound is \(H\le LI\), attained by the deterministic strategies above.
This also covers \(h=0\), the critical line, and both coefficient signs.

\paragraph{The quartic matrix.}
On \(c=d\), the Bell basis
\[
 \phi_\pm=(|00\rangle\pm|11\rangle)/\sqrt2,\qquad
 \psi_\pm=(|01\rangle\pm|10\rangle)/\sqrt2
\]
splits off \(\psi_-\), with eigenvalue \(-2hc\).
In the basis \((\phi_+,\phi_-,\psi_+)\), the remaining block is
\[
 H_3(c)=
 \begin{pmatrix}
 2hc&2a&0\\
 2a&2h(1+c-c^2)&2h(1-c)\sqrt{1-c^2}\\
 0&2h(1-c)\sqrt{1-c^2}&-2h(1+c-c^2)
 \end{pmatrix}.
\]
To remove the square root, use \(c=(1-t^2)/(1+t^2)\) and multiply
the slack $\lambda I_3-H_3(c)$ by \((1+t^2)^2\).
Writing
\[
 p_0=(1+t^2)^2,\qquad p_1=1-t^4,\qquad p_2=1+4t^2-t^4,
\]
the resulting quartic is
\begin{equation}
 K_{h,a,\lambda}(t)=
 \begin{pmatrix}
 \lambda p_0-2hp_1&-2ap_0&0\\
 -2ap_0&\lambda p_0-2hp_2&-8ht^3\\
 0&-8ht^3&\lambda p_0+2hp_2
 \end{pmatrix}.
 \label{eq:K}
\end{equation}
Finite \(t\ge0\) covers \(-1<c\le1\); \(c=-1\) follows by continuity.
Negative \(t\) adds no condition because
\(K(-t)=\operatorname{diag}(1,1,-1)K(t)\operatorname{diag}(1,1,-1)\).

Conversely, \(K(t)\succeq0\) for all real \(t\) forces the local bound.
Its upper-left block at \(t=0\) gives \(\lambda\ge2h+2|a|\);
the leading coefficient at \(t=\infty\) gives \(\lambda\ge-2h+2|a|\).
Thus \(\lambda\ge L\), which also bounds the omitted eigenvalue $-2hc$.
Equation~\eqref{eq:diagonal} applies for $|a|\le|h|$, and
Eq.~\eqref{eq:critical} gives the bound for $|a|\ge|h|$.
We have proved, for all real coefficients,
\begin{equation}
 (\lambda,a,h)\in{\cal C}
 \quad\Longleftrightarrow\quad
 K_{h,a,\lambda}(t)\succeq0\quad\text{for every }t\in\mathbb R .
 \label{eq:cone}
\end{equation}

\paragraph{A fixed affine Gram lift.}
Every real univariate matrix polynomial that is positive semidefinite (PSD)
for all real arguments can be written as a sum $\sum_j F_j(t)^TF_j(t)$
of polynomial matrix squares, including identically singular ones~\cite{hs}.
Since the diagonal entries of \(K\) have degree at most four, every square factor
has degree at most two: highest-degree real squares cannot cancel.
For $\xi=(\xi_1,\xi_2,\xi_3)^T\in\mathbb R^3$, a Gram representation
expresses $\xi^TK(t)\xi$ as a quadratic form in monomials in $t$ and $\xi$.
Average this representation under
$(t,\xi)\mapsto(-t,\operatorname{diag}(1,1,-1)\xi)$.
The even and odd monomials then form separate blocks, with vectors
\[
 Z_+=(\xi_1,\xi_2,t\xi_3,t^2\xi_1,t^2\xi_2)^T,\qquad
 Z_-=(\xi_3,t\xi_1,t\xi_2,t^2\xi_3)^T.
\]
Thus \(\xi^TK(t)\xi=Z_+^TGZ_++Z_-^TJZ_-\), with
\(G\succeq0\) of order five and \(J\succeq0\) of order four.
The sixteen coefficient equations are
\begin{equation}
\begin{aligned}
 g_{11}&=\lambda-2h,&g_{12}&=-2a,\\
 g_{22}&=\lambda-2h,&j_{11}&=\lambda+2h,\\
 g_{13}+j_{12}&=0,&g_{23}+j_{13}&=0,\\
 2g_{14}+j_{22}&=2\lambda,&g_{15}+g_{24}+j_{23}&=-4a,\\
 2g_{25}+j_{33}&=2\lambda-8h,&g_{33}+2j_{14}&=2\lambda+8h,\\
 g_{34}+j_{24}&=0,&g_{35}+j_{34}&=-8h,\\
 g_{44}&=\lambda+2h,&g_{45}&=-2a,\\
 g_{55}&=\lambda+2h,&j_{44}&=\lambda-2h .
\end{aligned}
\label{eq:rows}
\end{equation}
These equations include zero coefficients, lower polynomial degrees and every
rank-deficient case. Setting \(h=1\) describes all pairs $(a,\lambda)$ for
which $S+aM\le\lambda$ holds on $\W$. Restricting to \(0\le a\le1\)
requires the additional scalar inequalities $a\ge0$ and $a\le1$.

\section{Explicit membership matrices for the quantum region}

Introduce sixteen real variables \(z_1,\ldots,z_{16}\) and define
\begin{equation}
\begin{split}
 U(z)&=
 \begin{pmatrix}
 z_1&z_2/2&z_5/2&z_7&z_8/2\\
 z_2/2&z_3&z_6/2&z_8/2&z_9\\
 z_5/2&z_6/2&z_{10}&z_{11}/2&z_{12}/2\\
 z_7&z_8/2&z_{11}/2&z_{13}&z_{14}/2\\
 z_8/2&z_9&z_{12}/2&z_{14}/2&z_{15}
 \end{pmatrix},\\[2pt]
 V(z)&=
 \begin{pmatrix}
 z_4&z_5/2&z_6/2&z_{10}\\
 z_5/2&z_7&z_8/2&z_{11}/2\\
 z_6/2&z_8/2&z_9&z_{12}/2\\
 z_{10}&z_{11}/2&z_{12}/2&z_{16}
 \end{pmatrix}.
\end{split}\label{eq:dualmat}
\end{equation}
Membership is equivalent to feasibility of the following semidefinite
program (SDP):
\[
 (M,S)\in\W\quad\Longleftrightarrow\quad
 \text{there exists }z\in\mathbb R^{16}
 \text{ with }U(z)\succeq0,\quad V(z)\succeq0,
\]
and
\begin{align}
 1={}&z_1+z_3+z_4+2z_7+2z_9+2z_{10}+z_{13}+z_{15}+z_{16},\nonumber\\
 M={}&2z_2+4z_8+2z_{14},\nonumber\\
 S={}&2z_1+2z_3-2z_4+8z_9-8z_{10}+8z_{12}
       -2z_{13}-2z_{15}+2z_{16}.
 \label{eq:coords}
\end{align}

We first construct feasible $z$ from two-qubit states, proving that every
point in $\W$ passes this test.
Use $u,v\ge0$, $u^2+v^2=1$, and a real unit vector
$\xi=(\xi_1,\xi_2,\xi_3)$ in the triplet basis
$(\phi_+,\phi_-,\psi_+)$. The common local measurement angle has cosine $c=u^2-v^2$,
with sine $2uv$. Define
\[
 r_+=(u^2\xi_1,u^2\xi_2,uv\xi_3,v^2\xi_1,v^2\xi_2)^T,\qquad
 r_-=(u^2\xi_3,uv\xi_1,uv\xi_2,v^2\xi_3)^T.
\]
The matrices $U=r_+r_+^T$, $V=r_-r_-^T$ each have rank at most one and
have precisely the shared entries in Eq.~\eqref{eq:dualmat}. Explicitly, choose
\begin{equation}
\begin{split}
 z={}&(u^4\xi_1^2,\ 2u^4\xi_1\xi_2,\ u^4\xi_2^2,\ u^4\xi_3^2,\\
 &2u^3v\xi_1\xi_3,\ 2u^3v\xi_2\xi_3,\ u^2v^2\xi_1^2,
   \ 2u^2v^2\xi_1\xi_2,\\
 &u^2v^2\xi_2^2,\ u^2v^2\xi_3^2,\ 2uv^3\xi_1\xi_3,
   \ 2uv^3\xi_2\xi_3,\\
 &v^4\xi_1^2,\ 2v^4\xi_1\xi_2,\ v^4\xi_2^2,\ v^4\xi_3^2).
\end{split}
\label{eq:statewitness}
\end{equation}
Let $\ell(z)$ denote the linear form on the right-hand side of the
normalization equation in Eq.~\eqref{eq:coords}. Substitution gives $\ell(z)=1$ and
\begin{align*}
 M={}&4\xi_1\xi_2,\\
 S={}&2(u^4-v^4)\xi_1^2
 +2(u^4+4u^2v^2-v^4)(\xi_2^2-\xi_3^2)
 +16uv^3\xi_2\xi_3.
\end{align*}
These are the state's observed averages. The variables $u,v$ include
both endpoints without division or limits.

Let $\mathcal T$ be the set of pairs $(M,S)$ obtained from these real states
in the triplet subspace. It is compact and path-connected, being a continuous
image of the quarter-circle $u,v\ge0$, $u^2+v^2=1$, times the unit sphere $S^2$.
The singlet contributes only $(0,-2c)$, and
\[
 (0,-2c)=\frac{1+c}{2}(0,-2)+\frac{1-c}{2}(0,2).
\]
Both endpoints belong to $\mathcal T$: choose $\xi=(0,0,1)$ and
respectively $(u,v)=(1,0),(0,1)$. Since choices with $c=d$ attain every
support value, $\operatorname{conv}\mathcal T$ and $\W$ have the same
supporting inequalities. They are compact convex sets, hence
\begin{equation}
 \W=\operatorname{conv}\mathcal T.
 \label{eq:triplet-hull}
\end{equation}
Convex mixtures of the displayed witnesses therefore admit every point of
$\W$ into the SDP, including the boundary.

To prove that every feasible $z$ gives a point in $\W$, write Eq.~\eqref{eq:rows} as
$\mathcal A(G,J)=B(\lambda,a,h)$ in the displayed row order.
Here $\mathcal A$ maps the two Gram matrices to the sixteen left-hand sides,
and $B$ maps $(\lambda,a,h)$ to the corresponding right-hand sides.
With the trace inner product on symmetric matrices, the adjoint of
$\mathcal A$ satisfies $(U(z),V(z))=\mathcal A^*z$ and
$B^Tz=(\ell(z),-M(z),-S(z))$. Any feasible $z$ satisfies
\[
 \lambda-aM(z)-hS(z)=\tr[U(z)G]+\tr[V(z)J]\ge0
\]
for every valid supporting inequality, by Eq.~\eqref{eq:cone} and its Gram
representation. Thus its coordinates belong to $\W$. This proves the
membership test without invoking dual attainment.

\paragraph{Two strategies suffice.}
Every point of the convex hull of a compact connected set in the plane
lies on a segment joining two points of the set. Here is a direct proof.
Take $p\in\operatorname{conv}\mathcal T\setminus\mathcal T$ and project
$\mathcal T$ radially onto the unit circle about $p$. Its image is compact
and connected. If it contains no antipodal pair, it is an arc of length
less than $\pi$, contained in an open semicircle. The midpoint direction
then strictly separates $p$ from $\mathcal T$, a contradiction.
Thus $x-p=r n$ and $y-p=-s n$ for some $x,y\in\mathcal T$ and $r,s>0$,
and $p=(s x+r y)/(r+s)$.

Applied to Eq.~\eqref{eq:triplet-hull}, this shows that every $(M,S)\in\W$
is realized by a mixture of at most two real pure two-qubit strategies,
each with equal local measurement angles. Their witnesses give simultaneously
$\operatorname{rank}U,\operatorname{rank}V\le2$.
If the point is extreme, it belongs to $\mathcal T$ and both ranks can be
at most one. These rank bounds describe available witnesses; no rank
constraint is added to the convex SDP.

\paragraph{Compactness and the dual formulation.}
The normalized feasible auxiliary set is also compact. The following
argument gives an alternative dual proof of membership and can be useful
when optimizing over the lift.
Define the dual image
\[
 {\cal D}=\{B^Tz:U(z)\succeq0,\ V(z)\succeq0\}
\]
and note that it is closed. Indeed, the positive definite matrices
\[
 G_0=\operatorname{diag}(1,1,2,1,1),\qquad
 J_0=\operatorname{diag}(1,2,2,1)
\]
satisfy \(\mathcal A(G_0,J_0)=B(1,0,0)\), and hence
\(\ell(z)=\tr[U(z)G_0]+\tr[V(z)J_0]\).
If \(B^Tz\) converges, this weighted trace bounds both PSD matrices.
Every \(z_i\) occurs in their entries, so \(z\) is bounded and has a convergent
subsequence. The limit is feasible. Let ${\cal C}^*$ denote the dual cone,
whose elements have nonnegative inner product with every element of ${\cal C}$.
The strictly feasible pair \((G_0,J_0)\)
allows us to take the dual of the Gram representation, giving
\({\cal C}^*=\overline{\cal D}={\cal D}\).
Finally, the compact convex set \(\W\) is the intersection of its support
half-spaces, so \((M,S)\in\W\) exactly when \((1,-M,-S)\in{\cal C}^*\).
This proves the membership test. The normalized feasible \(z\)-set is itself
compact by the same trace argument.

\section{Minimal block order and the NPA comparison}

The set \(\W\) is compact and convex. The explicit lift also makes it
semialgebraic: its membership condition can be expressed by a finite
combination of polynomial equalities and inequalities.
Writing $\sxdeg(\W)$ for its semidefinite extension degree, Scheiderer's
planar theorem gives \(\sxdeg(\W)\le2\)~\cite{soc}.
A short argument shows that \(\W\) is not a polygon.

Let \(f(a)=\max_{\W}(S+aM)\), and \(w=2\sqrt2=f(0)\).
The identity
\[
 \mathrm{CHSH}(c,d)^2=4I+4rs\,Y\otimes Y
\]
shows that attaining \(w\) requires \(r=s=1\).
At \(c=d=0\), the top state is
\(v=\cos(\pi/8)\phi_-+\sin(\pi/8)\psi_+\), with both reduced states \(I/2\).
Thus every CHSH-maximizing correlation has \(M=0\).
For maximizers \(p_a\) of \(f(a)\), compactness then gives \(M(p_a)\to0\) as
\(a\downarrow0\), and
\[
 0\le\frac{f(a)-w}{a}\le M(p_a)\longrightarrow0.
\]
Yet \(f(a)>w\) for every \(a>0\): keeping \(c=d=0\), the normalized state
\((v+\varepsilon\phi_+)/\sqrt{1+\varepsilon^2}\) gives
\[
 \frac{w+4a\cos(\pi/8)\varepsilon}{1+\varepsilon^2}>w
 \quad\text{for}\quad \varepsilon=\frac{2a\cos(\pi/8)}w .
\]
A polygon would make \(f\) piecewise affine with finitely many pieces.
Its zero right derivative at zero would force an initial constant piece,
contradicting this strict increase. Therefore \(\W\) is nonpolyhedral, and
\[
 \boxed{\sxdeg(\W)=2.}
\]
The same proof gives extension degree two for
\(\{(a,\lambda):0\le a\le1,\ \lambda\ge f(a)\}\).

Write $\omega_{\mathrm{NPA}}^{(L)}(a)$ for the level-$L$ NPA upper bound on
$S+aM$. The prior NPA result applies to these support directions
\(h=1\)~\cite{prior}, with $T=1-a\downarrow0$:
\[
 f(1-T)=4-2T+\frac43T^3+O(T^4),\qquad
 \frac{\omega_{\mathrm{NPA}}^{(L)}(1-T)-f(1-T)}{T^3}\longrightarrow+\infty
\]
for every fixed finite level \(L\).
Consequently no fixed NPA projection onto \((M,S)\) equals \(\W\), while
Eqs.~\eqref{eq:dualmat}--\eqref{eq:coords} describe that same region exactly.
The nontermination theorem is proved in Part II. The finite lift and the
nonpolyhedrality argument above are independent of that proof.

\section{All signed asymmetric tilt branches}

The determinant minimization extends to arbitrary real marginal coefficients in
\[
 H_{\alpha,\beta}=\mathrm{CHSH}+\alpha A_0+\beta B_0.
\]
Put \(q=|\alpha|+|\beta|\). Its local bound is \(2+q\), attained for any tilt
signs. On $q=2$, combine the signed one-sided bounds from
Eq.~\eqref{eq:critical} with weights $|\alpha|/2$ and $|\beta|/2$.
This gives the same quantum bound. For $q>2$, first apply this result
to coefficients $2\alpha/q,2\beta/q$. The remaining marginal terms
are bounded above by $q-2$, giving $2+q$ again.

Suppose \(q<2\) and \(\lambda\ge2+q\). On a two-qubit Jordan block,
the coefficients of $u^2$ and $u$ in
$\det(uI+\lambda I-H_{\alpha,\beta})$ are
\[
 e_2=6\lambda^2-2\alpha^2-2\beta^2-8,\qquad
 e_3=4\lambda^3-4\lambda(\alpha^2+\beta^2+4)-8\alpha\beta\Omega .
\]
Both are nonnegative. For \(e_3\), a lower bound is
\(g(\lambda)=4\lambda^3-4\lambda(q^2+4)-4q^2\).
It increases on this range, and \(g(2+q)=32q+12q^2\ge0\).
Together with $e_1=4\lambda\ge0$, these inequalities show that the
determinant again decides positive semidefiniteness.

With \(x=(1-c)/2\), \(y=(1-d)/2\), direct expansion gives
\begin{align*}
 D(x,y)&=D_0+xy[C+64Ax+64By-256xy],\\
 A&=4-\alpha^2,\qquad B=4-\beta^2,\\
 C&=32\lambda\alpha\beta+64\alpha^2+64\beta^2-256,\\
 D_0&=[(\lambda-2)^2-(\alpha+\beta)^2]
       [(\lambda+2)^2-(\alpha-\beta)^2].
\end{align*}
Both \(A,B\) are positive. Exchange parties if necessary so that
\(|\alpha|\ge|\beta|\), and set \(\rho=A/B\in(0,1]\).
For a fixed product \(p=xy\), minimize \(Ax+Bp/x\) on \(p\le x\le1\).
For \(p\le\rho\), its minimizer lies on \(x=z,y=\rho z\).
For \(p>\rho\), it lies on \(x=1\). On that edge,
\[
 D(1,y)=-F_-(y)F_+(y),
\]
where
\[
 F_-=-(\lambda+\alpha)^2+\beta^2+4+4\beta(2y-1),\qquad
 F_+=(\lambda-\alpha)^2-\beta^2-4+4\beta(2y-1).
\]
Since \(\lambda\pm\alpha\ge2+|\beta|\), we have \(F_-\le0\le F_+\).
Thus it suffices to test the segment $(x,y)=(z,\rho z)$,
$0\le z\le1$, for every choice of tilt signs.

Along this segment, put
\(R=4z(1-z)\), \(T=4\rho z(1-\rho z)\), and \(\Omega=2-4\rho z^2\).
Congruence by \(\operatorname{diag}(\sqrt{RT},\sqrt R,\sqrt T,1)\) gives
\begin{equation}
 P(z)=
 \begin{pmatrix}
 RT(\lambda-\Omega-\alpha-\beta)&-2zRT&-2\rho zRT&RT\\
 -2zRT&R(\lambda+\Omega-\alpha+\beta)&RT&2\rho zR\\
 -2\rho zRT&RT&T(\lambda+\Omega+\alpha-\beta)&2zT\\
 RT&2\rho zR&2zT&\lambda-\Omega+\alpha+\beta
 \end{pmatrix}.
 \label{eq:rayP}
\end{equation}
This matrix has degree at most six. The congruence is invertible on \(0<z<1\);
the original slack at both endpoints follows by continuity from that interval.
The matrix interval sum-of-squares theorem gives
\begin{equation}
 P(z)=Y_3(z)^TUY_3(z)+z(1-z)Y_2(z)^TVY_2(z),\qquad
 U\succeq0,\quad V\succeq0,
 \label{eq:interval}
\end{equation}
with \(Y_j=(I_4,zI_4,\ldots,z^jI_4)^T\), and block orders \(16,12\).
One can obtain this interval theorem directly from the real-line matrix theorem:
apply it to \((1+t^2)^6P(t^2/(1+t^2))\), split its factors into even and odd parts,
and divide back. Even factors become polynomials of degree at most three in \(z\);
odd factors give \(\sqrt{z(1-z)}\) times polynomials of degree at most two.
This argument also includes singular matrices and endpoint zeros.

For fixed \(\alpha,\beta\), coefficient matching in Eq.~\eqref{eq:interval}
is affine in \(\lambda,U,V\). Include \(\lambda\ge2+q\) explicitly.
This yields an exact SDP of fixed size whenever
$|\alpha|+|\beta|<2$, including either zero-tilt axis and CHSH.
For $|\alpha|+|\beta|\ge2$, the local-bound certificates above apply.
The ratio \(\rho\) still depends nonlinearly on the two independent tilts:
this construction computes their support values, while the common-tilt
construction above gives an affine membership lift.


\section{Thirty-two exact projections and their NPA intersections}
\label{finite:sec:orbit}

Write $a_x=\langle A_x\rangle$, $b_y=\langle B_y\rangle$, and
$E_{ij}=\langle A_iB_j\rangle$. For a sign pattern $\chi\in\{\pm1\}^{2\times2}$ satisfying
$\prod_{ij}\chi_{ij}=-1$, and $x,y\in\{0,1\}$, set
\[
 M_{\chi,xy}=a_x+\chi_{xy}b_y,\qquad
 S_\chi=\sum_{ij}\chi_{ij}E_{ij}.
\]
The quantum image of $(M_{\chi,xy},S_\chi)$ is exactly $\W$.
Here is an explicit physical relabelling. In the canonical observables
replace
\[
 A_0'=A_x,\quad A_1'=\varepsilon_1 A_{1-x},\qquad
 B_0'=\eta_0B_y,\quad B_1'=\eta_1B_{1-y},
\]
where
\[
 \eta_0=\chi_{xy},\qquad
 \eta_1=\chi_{x,1-y},\qquad
 \varepsilon_1=\chi_{1-x,y}/\eta_0.
\]
Then $a'_0+b'_0=M_{\chi,xy}$. The first three CHSH coefficients
are the corresponding entries of $\chi$, and the fourth is
$-\varepsilon_1\eta_1=\chi_{1-x,1-y}$ because the product of
the four entries is $-1$. The relabelling is invertible and preserves
quantum realizability, which proves equality of the images.

These are thirty-two distinct two-dimensional subspaces of the
space of Bell functionals. The marginal part identifies $x,y$.
Equality of the correlation directions would require
$\chi'=\chi$ or $\chi'=-\chi$.
The latter changes the relative sign of the two nonzero marginal
coefficients, so the marginal directions are not proportional.
Thus only the identical pair gives the same plane.

Let $\Q_L$ be a fixed standard NPA relaxation and define
\[
 \mathcal R_L=\Q_L\cap
 \bigcap_{\chi,x,y}\{p:\pi_{\chi,xy}(p)\in\W\},\qquad
 \pi_{\chi,xy}(p)=(M_{\chi,xy},S_\chi).
\]
Every quantum correlation satisfies these constraints, and each projection
of $\mathcal R_L$ lies in $\W$:
\[
 \Q_{222}\subseteq\mathcal R_L\subseteq
 \pi_{\chi,xy}^{-1}(\W).
\]
For any real $h,\alpha$, these two inclusions imply
\[
 \sup_{p\in\mathcal R_L}
 [hS_\chi(p)+\alpha M_{\chi,xy}(p)]
 =
 \sup_{p\in\Q_{222}}
 [hS_\chi(p)+\alpha M_{\chi,xy}(p)] .
\]
Part II gives directions on the canonical plane for which the
level-$L$ support is strictly larger. Hence
$\mathcal R_L\subsetneq\Q_L$ for every finite $L$.

The rational data file
\texttt{ALL32\_EXACT\_PROJECTION\_CONSTRAINTS.json}
implements each copy of Eqs.~\eqref{eq:dualmat}--\eqref{eq:coords}
with the two appropriate linear forms. The added constraints use
64 PSD blocks, alternating orders five and four, total order 288,
512 auxiliary variables and 96 affine equations. The generator enumerates
physical relabellings and checks the distinct maps exactly.
The inclusions above establish exactness on the stated planes;
they do not establish quantum realizability of a full correlation whose
projections satisfy them.

\clearpage
\part*{II. The complete finite-NPA obstruction}
\setcounter{section}{0}\setcounter{equation}{0}\setcounter{theorem}{0}
\renewcommand{\thesection}{S\arabic{section}}
\renewcommand{\thesubsection}{\thesection.\arabic{subsection}}
\renewcommand{\theequation}{S\arabic{equation}}
\renewcommand{\thetheorem}{S\arabic{theorem}}
\renewcommand{\theHequation}{S.\arabic{equation}}
\renewcommand{\theHsection}{S.\arabic{section}}
% Replacement body for Part II, after its heading and counter/numbering setup.
\section{Bell values, quantum models, and finite moment certificates}
\label{supp:sec:bell}\label{supp:sec:jordan}\label{supp:sec:npa}
Alice and Bob choose $x,y\in\{0,1\}$ and obtain $a,b\in\{\pm1\}$.
A valid correlation is a nonnegative, normalized, no-signaling table
$p(a,b|x,y)$. Its four marginal expectations and four correlators determine it:
\begin{equation}
p(a,b|x,y)=\tfrac14[1+a\langle A_x\rangle+b\langle B_y\rangle
+ab\langle A_xB_y\rangle].
\label{supp:eq:behavior-reconstruction}
\end{equation}
Here the expectations are the averages of $a$, $b$, and $ab$ under the
corresponding marginals or joint distribution. For $0\leq\alpha<1$, set
\begin{equation}
h_\alpha=A_0B_0+A_0B_1+A_1B_0-A_1B_1+\alpha(A_0+B_0),
\qquad \beta_\alpha(p)=\langle h_\alpha\rangle_p.
\label{supp:eq:bell-functional}
\end{equation}
The CHSH bound $2$ and the two marginal bounds give the local maximum
$\omega_{\rm L}(\alpha)=2+2\alpha$, attained when all outputs are $+1$.
In this part we abbreviate $\Q_{222}$ to $\Q$.
Write $\omega_{\rm Q}(\alpha)=\max_{p\in\Q}\beta_\alpha(p)$ and
$\Delta_\alpha=\omega_{\rm Q}(\alpha)-\omega_{\rm L}(\alpha)$,
the quantum violation above the local bound.

The quantum set $\Q$ consists of
$p(a,b|x,y)=\Tr[\rho(E_{a|x}\otimes F_{b|y})]$, with a density operator
$\rho$ and local binary positive-operator-valued measurements (POVMs),
$\{E_{a|x}\}_a$ and $\{F_{b|y}\}_b$. A common local Naimark dilation preserves
all probabilities and makes both settings projective~\cite{Naimark1940};
the commuting-model version follows from Proposition~3.4(b) of
Ref.~\cite{Fritz2012}. The differences of the two outcome projectors
are self-adjoint reflections,
$A_x^2=B_y^2=I$, with projectors $(I+aA_x)/2$, $(I+bB_y)/2$.

Let $\Q_{\rm fd}$ denote the finite-dimensional tensor-product quantum set.
Jordan's lemma decomposes each finite-dimensional local pair into invariant
blocks of dimension at most two~\cite{Halmos1969,Masanes2005,PanwarPandyaWiesniak2022}.
For block projections $P_{ij}=\Pi_i^A\otimes\Pi_j^B$, the weights
$r_{ij}=\Tr(P_{ij}\rho)$ and states $P_{ij}\rho P_{ij}/r_{ij}$ give
$p=\sum_{ij:r_{ij}>0}r_{ij}p_{ij}$. Conversely, correlated classical labels
realize any finite mixture by local direct sums. If $\mathcal K_2$ denotes
the two-qubit correlations with projective measurements, including the
projectors $0$ and $I$, then $\mathcal Q_{\rm fd}=\operatorname{conv}\mathcal K_2$.
Carath\'eodory's theorem requires at most nine components in the
eight-dimensional correlation space. The state and projector parameter sets
are compact, so the mixture set is the compact continuous image of
$\Delta_8\times\mathcal K_2^9$, where $\Delta_8$ is the simplex of nine
nonnegative weights summing to one. Finite-dimensional approximation adds no
tensor-product correlations. In the commuting model, denoted by $\Q_{\rm c}$,
Alice's and Bob's operators act on one Hilbert space and commute with each other.
The local two-projection algebra is nuclear
\cite[Sec.~4, Example~7(2)]{ScholzWerner2008}; Proposition~3.4 and
Remark~3.6 of Ref.~\cite{Fritz2012} therefore identify the commuting set with
the tensor-product closure. Thus, with complex Hilbert spaces throughout,
\begin{equation}
\Q_{\rm fd}=\Q=\overline{\Q_{\rm fd}}=\Q_{\rm c}
=\operatorname{conv}\mathcal K_2.
\label{supp:eq:quantum-model-equality}
\end{equation}
In particular, NPA converges to $\Q$ here~\cite[Theorem~8]{NPA2008}.

To define the finite NPA levels, we use formal observable products, called
words. The group relations express that the observables square to the
identity and that Alice's observables commute with Bob's. With $e$ denoting
the identity, the group is
\begin{equation}
G=\langle a_0,a_1,b_0,b_1:a_x^2=b_y^2=e,\ a_xb_y=b_ya_x\rangle.
\label{supp:eq:group-presentation}
\end{equation}
Its complex group algebra consists of finite complex linear combinations
of these words. Reduced words have Alice's alternating letters before Bob's
alternating letters; the adjoint reverses a word and conjugates its coefficient. Let
$\mathcal W_L$ contain the reduced words of length at most $L$, let
$V_L=\Span_\C\mathcal W_L$, and enumerate its words as $\bm w$.
A candidate moment assignment associates a complex number $\ell(w)$ with
each word through length $2L$, extends linearly, and is Hermitian when
$\ell(w^*)=\overline{\ell(w)}$. It need not come from a physical realization.
Its moment matrix and the resulting NPA bound are
\begin{equation}
\Gamma_L(\ell)=(\ell(u^*v))_{u,v\in\mathcal W_L},\qquad
\omega_{\rm NPA}^{(L)}(\alpha)
=\max_{\ell(e)=1,\ \Gamma_L(\ell)\succeq0}\ell(h_\alpha),
\label{supp:eq:primal-NPA}
\end{equation}
where $h_\alpha$ now uses the abstract generators. For a measurement
representation $\pi$ and a unit state vector $|\psi\rangle$, physical moments
$\ell(w)=\langle\psi|\pi(w)|\psi\rangle$ make $\Gamma_L$ a Gram matrix;
therefore $\omega_{\rm NPA}^{(L)}\geq\omega_{\rm Q}$.
We call $\omega_{\rm NPA}^{(L)}-\omega_{\rm Q}$ the NPA approximation error.
Level $1$ is an observable moment bound and need not give nonnegative
probabilities. We extract correlations only at levels $d\geq2$ below.

An identity $\lambda e-h_\alpha=\sum_jp_j^*p_j$ in the group algebra,
with $p_j\in V_L$, is an upper-bound certificate: evaluation on any
measurements gives $\lambda I-\pi(h_\alpha)\succeq0$. Its Gram matrix is the
variable in the dual problem below.

\begin{proposition}[Attainment and Gram trace bound]
\label{supp:prop:duality}\label{supp:lem:trace-identity}
For every finite $L\geq1$ and $0\leq\alpha\leq1$, the primal maximum in
Eq.~\eqref{supp:eq:primal-NPA} and the dual minimum below are attained:
\begin{equation}
\omega_{\rm NPA}^{(L)}(\alpha)
=\min\{\lambda\in\R:\lambda e-h_\alpha=\bm w^*Q\bm w,\ Q\succeq0\}
\label{supp:eq:dual-NPA}
\end{equation}
The Gram matrix may be real symmetric. Every such certificate satisfies
the following bounds, where matrix norms are operator norms:
\begin{equation}
\Tr Q=\lambda,\qquad \|Q\|\leq\lambda,\qquad
\|Q^{1/2}\|\leq\sqrt\lambda.
\label{supp:eq:trace-Q}
\end{equation}
\end{proposition}
\begin{proof}
The canonical trace $\tau(\sum_gc_gg)=c_e$ obeys
$\tau(u^*v)=\delta_{uv}$, so $\Gamma_L(\tau)=I\succ0$.
For feasible $\ell$, the two-by-two principal minors give
$|\ell(u^*v)|^2\leq\ell(u^*u)\ell(v^*v)=1$.
Every word through length $2L$ is $u^*v$ with $u,v\in\mathcal W_L$
(split the word into two pieces of length at most $L$ and invert the first).
Thus the feasible set is closed and bounded in a finite-dimensional space,
so the maximum is attained. Strict feasibility and finiteness give dual
attainment and equality by Slater's theorem~\cite{BoydVandenberghe2004}.
The dual identity is precisely Eq.~\eqref{supp:eq:dual-NPA}: writing
$Q=C^*C$ gives a sum $\sum_jp_j^*p_j$ with $p_j\in V_L$, and conversely
every such sum gives a positive Gram matrix. The multiplication table and
$h_\alpha$ have real coefficients, so averaging $Q$ with $\overline Q$
preserves the identity and positivity. Applying $\tau$ gives
$\lambda=\sum_{ij}Q_{ij}\tau(w_i^*w_j)=\Tr Q$, since $\tau(h_\alpha)=0$;
positivity gives the norm bounds. The moment bounds also imply
$\omega_{\rm NPA}^{(L)}(\alpha)\leq4+2\alpha\leq6$.
\end{proof}

For comparison, $\mathcal W_1=\{e,a_0,a_1,b_0,b_1\}$ already recovers
the CHSH maximum $2\sqrt2$~\cite{Tsirelson1980,NPA2007}. The list
$\mathcal W_{1+AB}=\mathcal W_1\cup\{a_xb_y\}_{x,y}$ gives tight
certificates for one-sided tilted CHSH,
$h_0+\gamma a_0$ for $0\leq\gamma<2$, and self-tests its maximizers
\cite{AMP2012,BampsPironio2015,SupicBowles2020}. Thus level $1+AB$ is
exact when either marginal tilt vanishes~\cite{GigenaEtAl2025}.

\section{The cubic quantum violation}
\label{supp:sec:endpoint}\label{supp:sec:cubic-scale}
On a nontrivial Jordan block local basis choices give
$A_0=B_0=Z$, $A_1=cZ+rX$, $B_1=dZ+sX$, where
$c,d\in[-1,1]$, $r=\sqrt{1-c^2}$, and $s=\sqrt{1-d^2}$.
The endpoints include one-dimensional blocks: the assignment
$(A_0,A_1)=(\epsilon_0,\epsilon_1)$ occurs on the $Z$ eigenspace
$\epsilon_0$ at $c=\epsilon_0\epsilon_1$, and similarly for Bob.
Put $\alpha=1-T$, $\lambda_0(T)=4-2T$, $u=1-c$, $v=1-d$, and
$\Omega=2-uv$. The evaluated Bell operator is
\begin{equation}
H_T(c,d)=\begin{pmatrix}
\Omega+2(1-T)&us&vr&-rs\\
us&-\Omega&-rs&-vr\\
vr&-rs&-\Omega&-us\\
-rs&-vr&-us&\Omega-2(1-T)
\end{pmatrix}.
\label{supp:eq:block-matrix}
\end{equation}
Equation~\eqref{supp:eq:quantum-model-equality} gives
$\omega_{\rm Q}(1-T)=\max_{c,d\in[-1,1]}\lambda_{\max}H_T(c,d)$,
attained on two qubits. For later use, multiplication by a real vector
$\eta=(1,b_1,b_2,b_3)^T$ yields the exact identity, also valid for signed
$r,s$ with $r^2=1-c^2$, $s^2=1-d^2$,
\begin{align}
\mathcal E_T(\eta):={}&\langle\eta,(H_T-\lambda_0(T)I)\eta\rangle\nonumber\\
={}&-uv+2us b_1+2vr b_2-2rs b_3
+(-6+2T+uv)(b_1^2+b_2^2)\nonumber\\
&+(-4+4T-uv)b_3^2-2rs b_1b_2-2vr b_1b_3-2us b_2b_3.
\label{supp:eq:excess-quadratic-form}
\end{align}
The normalized vector's Bell excess is $\mathcal E_T(\eta)/\|\eta\|^2$.

Theorem~1 and Eqs.~(18)--(19) of Ref.~\cite{GigenaEtAl2025} identify
$\omega_{\rm Q}(\alpha)$ as the largest real root in $\lambda$ of
\begin{align}
F(\lambda,\alpha)={}&\lambda^4+(4-\alpha^2)\lambda^3
+(\tfrac{11}{4}\alpha^4-12\alpha^2-4)\lambda^2\nonumber\\
&+(2\alpha^6-\alpha^4-20\alpha^2-32)\lambda
+5\alpha^6-21\alpha^4+16\alpha^2-32,
\label{supp:eq:exact-quartic}
\end{align}
with common optimal cosine
\begin{equation}
c_*(\alpha)=\tfrac18[3\alpha^2-4+
\sqrt{16+9\alpha^4+8\alpha^2(2\omega_{\rm Q}(\alpha)-1)}].
\label{supp:eq:exact-optimal-cosine}
\end{equation}
At $\alpha=1$,
$F(\lambda,1)=(\lambda-4)(4\lambda^3+28\lambda^2+59\lambda+32)/4$.
The cubic is positive for $\lambda\geq4$. The local lower bound and
$\omega_{\rm Q}(\alpha)\leq4+2\alpha$ force every endpoint limit to be
the root $4$. Changing the tilt by $T$ changes any Bell value by at most
$2T$, so $\omega_{\rm Q}(1-T)=4+O(T)$. Since
\[
F(2+2\alpha,\alpha)=4(\alpha-1)^3(\alpha+2)^4,
\qquad \partial_\lambda F(4,1)=243,
\]
the mean-value theorem at fixed $\alpha=1-T$ gives
\begin{equation}
\Delta_{1-T}=\frac{4T^3(3-T)^4}
{\partial_\lambda F(\widetilde\lambda_T,1-T)}
=\frac{324}{243}T^3+O(T^4)=\frac43T^3+O(T^4),
\label{supp:eq:quantum-advantage}
\end{equation}
where $\widetilde\lambda_T$ lies between the local and quantum values.
The nonzero derivative also gives a unique analytic root near $4$.
Writing this root as
$4-2T+\frac43T^3+c_4T^4+c_5T^5+O(T^6)$ and substituting into $F=0$ gives
$27(9c_4+4)T^4+(243c_5-243c_4-120)T^5+O(T^6)=0$.
Consequently
\begin{align}
\omega_{\rm Q}(1-T)&=4-2T+\tfrac43T^3-\tfrac49T^4
+\tfrac4{81}T^5+O(T^6),
\label{supp:eq:exact-endpoint-expansion}\\
c_*(1-T)&=1-2T+\tfrac89T^2+\tfrac{20}{81}T^3+O(T^4).
\label{supp:eq:optimal-cosine-expansion}
\end{align}
Analyticity gives constants $C,\delta>0$ bounding each remainder by $CT^k$
at its displayed order for $0\leq T<\delta$. In particular the quartic
correction to the quantum value divided by $T^4$ tends to $-4/9$.
Since $1-\cos\theta=\theta^2/2+O(\theta^4)$, the optimal angles obey
$\theta_*\sim2\sqrt T$; the cited self-testing theorem gives this scale
for every maximizer in the canonical form.

This endpoint has a detector-loss interpretation~\cite{MiklinEtAl2022,GigenaEtAl2025}.
Let each detector fire independently with common probability $\eta$,
independently also of settings and ideal outcomes, and assign $+1$ on
no-detection. Then
\begin{align*}
\langle A_xB_y\rangle_\eta
&=\eta^2\langle A_xB_y\rangle+
\eta(1-\eta)(\langle A_x\rangle+\langle B_y\rangle)+(1-\eta)^2,\\
\mathrm{CHSH}_\eta(p)&=\eta^2\beta_\alpha(p)+2(1-\eta)^2,
\qquad \alpha=2(1-\eta)/\eta.
\end{align*}
The setting-$1$ marginals cancel in CHSH. The local bound becomes $2$,
and both the quantum violation and the NPA error are multiplied by
$\eta^2$, preserving their ratio. Here $\alpha\uparrow1$ corresponds
to $\eta\downarrow2/3$~\cite{Eberhard1993}; since
$T=(3\eta-2)/\eta=(9/2)(\eta-2/3)+O((\eta-2/3)^2)$,
\begin{equation}
\eta^2\Delta_\alpha=54(\eta-2/3)^3+O((\eta-2/3)^4).
\label{supp:eq:loss-cubic}
\end{equation}

\section{A two-parameter trial state and the Motzkin polynomial}
\label{supp:sec:motzkin}\label{supp:sec:not-sos}
The quantum violation is of order $T^3$, and the optimal measurement angles
are of order $\sqrt T$. Set $T=t^2$ so that these scales become $t^6$
and $t$, and vary the two angles independently:
\begin{equation}
A_0=B_0=Z,\quad A_1=\cos(tx)Z+\sin(tx)X,\quad
B_1=\cos(ty)Z+\sin(ty)X,\qquad x,y\in\R.
\label{supp:eq:two-parameter-observables}
\end{equation}
Write $H(t;x,y)$ for the evaluation of $h_{1-t^2}$ on these observables.
A certificate must hold for every $(x,y)$, although the canonical rescaled
optimal angles approach $(x,y)=(2,2)$. In what follows
$O_K(t^k)$ bounds the scalar, Euclidean, or operator norm uniformly on each
compact $K\subset\R^2$ by $C_K|t|^k$ for $0<|t|<\delta_K$.

At $t=0$, $H=\operatorname{diag}(4,-2,-2,0)$; the gaps from $|00\rangle$
are $6,6,4$. Its couplings to the other basis vectors are
$x^2yt^3/2+O_K(t^5)$, $xy^2t^3/2+O_K(t^5)$, and
$-xyt^2+O_K(t^4)$. Take
$\zeta_t=(1,At^3,Bt^3,Ct^2)^T$ and
$\Lambda_t=4-2t^2+4t^6/3$.
We choose $A,B,C$ to make the expectation of $\Lambda_tI-H(t)$ vanish
below order $t^6$.
Equation~\eqref{supp:eq:excess-quadratic-form} gives the leading gap
$(4C+xy)^2t^4/4$. Hence choose $C=-xy/4$. The next coefficient is
\begin{align}
\langle\zeta_t,(\Lambda_tI-H(t))\zeta_t\rangle
=t^6\biggl[&6(A-x^2y/12)^2+6(B-xy^2/12)^2\nonumber\\
&+\tfrac43+\frac{x^4y^2+x^2y^4}{48}-\frac{x^2y^2}{4}\biggr]
+O_K(t^8).
\label{supp:eq:gap-completed-squares}
\end{align}
Indeed, before completing the two squares their contribution is
$6A^2-Ax^2y+6B^2-Bxy^2$, and the coefficient of $x^4y^2+x^2y^4$
is $1/16$; completing the squares subtracts $1/24$, leaving $1/48$.
Choose $A=x^2y/12$ and $B=xy^2/12$. Define the resulting polynomial
vector $\xi_t$ and the Motzkin polynomial $M(X,Y)$ by
\begin{equation}
\xi_t=(1,t^3x^2y/12,t^3xy^2/12,-t^2xy/4)^T,
\qquad
M(X,Y)=X^4Y^2+X^2Y^4-3X^2Y^2+1
\label{supp:eq:scaled-test-vector}
\end{equation}
Then $\xi_t$ satisfies
\begin{equation}
\langle\xi_t,[\omega_{\rm Q}(1-t^2)I-H(t;x,y)]\xi_t\rangle
=\tfrac43t^6M(x/2,y/2)+O_K(t^8).
\label{supp:eq:motzkin-limit}
\end{equation}
To check the remainder, sine and cosine have uniform Taylor remainders on
$K$. Conjugation by $U=Z\otimes Z$ gives
$H(-t)=UH(t)U$ and $\xi_{-t}=U\xi_t$, so the expectation is even.
Replacing $\Lambda_t$ by the exact quantum value changes it by
$[-4t^8/9+4t^{10}/81+O(t^{12})]\|\xi_t\|^2=O_K(t^8)$, since
\begin{equation}
\|\xi_t\|^2=1+x^2y^2t^4/16+(x^4y^2+x^2y^4)t^6/144.
\label{supp:eq:test-vector-norm}
\end{equation}
Equivalently, $\xi_t$ is the top eigenvector, with first component one,
through order three: its leading component equations are
$6A=x^2y/2$, $6B=xy^2/2$, $4C=-xy$. The simple eigenvalue at $t=0$
and its fixed spectral gap give an $O_K(t^4)$ vector remainder.

The normalized state $\psi_t=\xi_t/\|\xi_t\|$ has Bell value
$b_t=\langle\psi_t,H(t)\psi_t\rangle$, and
\begin{equation}
\frac{\omega_{\rm Q}(1-t^2)-b_t(x,y)}{\Delta_{1-t^2}}
\longrightarrow M(x/2,y/2)
\label{supp:eq:normalized-bell-gap}
\end{equation}
locally uniformly. On either axis $\xi_t=|00\rangle$ is an eigenvector
at the local bound, so this ratio is exactly $1$. At $(x,y)=(2,2)$ it
tends to zero. The gap in this formula is a physical strategy's gap from
the quantum bound; the NPA error will add a separate nonnegative term.

\begin{lemma}[Shifted Motzkin polynomial]
\label{supp:lem:motzkin-not-sos}
$M$ is nonnegative, but $M+c$ is not a sum of real polynomial squares
for any $c\geq0$.
\end{lemma}
\begin{proof}
Arithmetic--geometric mean gives $X^4Y^2+X^2Y^4+1\geq3X^2Y^2$.
If $M+c=\sum_jq_j^2$, the highest homogeneous terms cannot cancel,
so $\deg q_j\leq3$. On either axis the sum is the constant $1+c$;
each restricted polynomial is bounded and therefore constant. Thus
$q_j=a_j+b_jXY+c_jX^2Y+d_jXY^2$. The $X^2Y^2$ coefficient of
$\sum_jq_j^2$ is $\sum_jb_j^2\geq0$, contradicting $-3$.
\end{proof}
An alternative proof uses the Newton polytope $\Newt(p)$, the convex hull
of the exponent vectors of monomials with nonzero coefficients in $p$.
The usual half-Newton-polytope property gives the same monomial restriction:
if $p=\sum_jq_j^2$, then
$\Newt(q_j)\subseteq\frac12\Newt(p)$~\cite[Theorem~1]{Reznick1978}.
The homogeneous statement gives the polynomial statement by homogenization
and then setting the extra variable to one. Here
$\Newt(M+c)=\operatorname{conv}\{(0,0),(4,2),(2,4)\}$, and its half
has integer points $(0,0),(1,1),(2,1),(1,2)$, exactly the four monomials
above. Independent angles are essential: the restriction
$M(X,X)=(X^2-1)^2+2[X(X^2-1)]^2$, even after adding $c\geq0$, is a
sum of squares.

\section{Polynomial limits of fixed certificates and divergence}
\label{supp:sec:compactness}\label{supp:sec:main-proof}
An evaluated certificate gives squared norms of analytic functions. To
obtain polynomial square factors in its limit, we first control the Taylor
remainder and then bound coefficients in one fixed polynomial space.
For a fixed list $\mathcal W_L=(w_1,\ldots,w_N)$, let
\begin{equation}
z_t(x,y)=(\pi_{t,x,y}(w_i)\xi_t(x,y))_{i=1}^N\in\C^{4N},
\label{supp:eq:z-vector}
\end{equation}
where $\pi$ evaluates Eq.~\eqref{supp:eq:two-parameter-observables}.

\begin{lemma}[Fixed-level sum-of-squares limit]
\label{supp:lem:fixed-level-limit}
Let $t_n\downarrow0$, $Q_n\succeq0$ be $N\times N$ matrices with
$\sup_n\Tr Q_n<\infty$, and $m\geq0$ be an integer. If
\begin{equation}
t_n^{-2m}\|(Q_n^{1/2}\otimes I_4)z_{t_n}(x,y)\|^2\longrightarrow P(x,y)
\label{supp:eq:compactness-assumption}
\end{equation}
locally uniformly for a real polynomial $P$, then $P$ is a sum of squares
of real polynomials.
\end{lemma}
\begin{proof}
A sine or cosine coefficient of order $k$ in $t$ is a polynomial in $x,y$
of total degree at most $k$; the entries of $\xi_t$ have the same property. Multiplication
adds these orders. Since the word list is fixed and finite, Taylor's
theorem gives, on every compact $K$,
\begin{equation}
z_t=\sum_{r=0}^m t^rz_r+R_{m+1},\qquad
\deg z_r\leq r,\qquad \sup_K\|R_{m+1}\|\leq C_K|t|^{m+1}.
\label{supp:eq:z-Taylor}
\end{equation}
Put $C_n=Q_n^{1/2}$, keeping all $N$ rows even if its rank changes.
The trace bound gives $\|C_n\|\leq\sqrt{\sup_n\Tr Q_n}$.
Define the polynomial vector $p_n$ of degree at most $m$ and the
analytic vector $v_n$ by
\begin{equation}
p_n=t_n^{-m}(C_n\otimes I_4)\sum_{r=0}^m t_n^rz_r,
\qquad v_n=t_n^{-m}(C_n\otimes I_4)z_{t_n},
\qquad \sup_K\|p_n-v_n\|=O_K(t_n),
\label{supp:eq:pn-definition}
\end{equation}
The vectors $p_n$ are uniformly bounded on $K_0=[-1,1]^2$, since $\|v_n\|^2\to P$
uniformly there. To bound coefficients despite the negative powers of
$t_n$, choose $m+1$ fixed distinct nodes $s_i\in[-1,1]$. Their fixed
Lagrange polynomials $L_i$ give, for each coordinate $p$ of $p_n$,
\[
p(x,y)=\sum_{i,j=0}^m p(s_i,s_j)L_i(x)L_j(y).
\]
The bounded grid values therefore bound every coefficient of $p_n$.
One subsequence makes all finitely many complex coefficients converge,
giving a polynomial vector $q$ and locally uniform convergence $p_n\to q$.
On $K_0$,
$|\|p_n\|^2-\|v_n\|^2|\leq(\|p_n\|+\|v_n\|)\|p_n-v_n\|\to0$,
so $P=\|q\|^2$ there. This is a polynomial identity on an open set and
thus on $\R^2$. Splitting each coordinate of $q$ into real and imaginary
parts expresses $P$ as real polynomial squares.
\end{proof}

\begin{theorem}
\label{supp:thm:fixed-level}
For every fixed finite $L\geq1$,
\begin{equation}
\lim_{\alpha\uparrow1}
\frac{\omega_{\rm NPA}^{(L)}(\alpha)-\omega_{\rm Q}(\alpha)}{(1-\alpha)^3}
=+\infty,
\qquad
\lim_{\alpha\uparrow1}
\frac{\omega_{\rm NPA}^{(L)}(\alpha)-\omega_{\rm Q}(\alpha)}{\Delta_\alpha}
=+\infty.
\label{supp:eq:main-divergence}
\end{equation}
\end{theorem}
\begin{proof}
If the first divergence failed, nonnegativity of the error would give
$t_n\downarrow0$ with
$\omega_{\rm NPA}^{(L)}(1-t_n^2)=\omega_{\rm Q}(1-t_n^2)+r_nt_n^6$,
where, after a subsequence, $r_n\to r\geq0$.
Proposition~\ref{supp:prop:duality} gives exact Gram certificates
\[
\omega_{\rm NPA}^{(L)}(1-t_n^2)e-h_{1-t_n^2}=\bm w^*Q_n\bm w,
\qquad \Tr Q_n=\omega_{\rm NPA}^{(L)}(1-t_n^2)\leq6.
\]
Evaluate on the observables and $\xi_{t_n}$, then split the expectation
into the quantum strategy gap and $r_nt_n^6\|\xi_{t_n}\|^2$.
Equations~\eqref{supp:eq:motzkin-limit} and \eqref{supp:eq:test-vector-norm}
give, locally uniformly,
\begin{equation}
t_n^{-6}\|(Q_n^{1/2}\otimes I_4)z_{t_n}\|^2
\longrightarrow\tfrac43M(x/2,y/2)+r.
\label{supp:eq:certificate-limit}
\end{equation}
Lemma~\ref{supp:lem:fixed-level-limit} with $m=3$ makes this polynomial
a sum of real squares. Substituting $(x,y)=(2X,2Y)$ and multiplying by
$3/4$ would express $M(X,Y)+3r/4$ as such a sum, contrary to
Lemma~\ref{supp:lem:motzkin-not-sos}. The second limit follows from
$\Delta_\alpha/(1-\alpha)^3\to4/3$.
\end{proof}
The absolute NPA error remains bounded. The divergence concerns its ratio
to a quantum violation that vanishes cubically. For each fixed $L$ it
also gives a full interval $1-\delta_L<\alpha<1$ with a strict gap;
$\delta_L$ may depend on $L$. No assertion that one fixed
$\alpha<1$ is missed at every finite level is needed.

\section{Valid correlations, arbitrary finite word lists, and a local limit}
\label{supp:sec:transfer}
Fix any finite list $\mathcal S=(S_1,\ldots,S_k)$ of words in the
measurement projectors. Define $\mathcal N_{\mathcal S}$ to contain the
valid no-signaling correlations for which a normalized Hermitian functional
$\ell$ satisfies
\begin{equation}
\Gamma_{\mathcal S}=(\ell(S_i^*S_j))_{ij}\succeq0,\qquad
p(a,b|x,y)=\ell(E_{a|x}F_{b|y}),
\label{supp:eq:selected-moment-matrix}
\end{equation}
together with the identities generated by
$E_{a|x}^*=E_{a|x}$,
$E_{a|x}E_{a'|x}=\delta_{aa'}E_{a|x}$, $\sum_aE_{a|x}=I$,
the analogous Bob relations, and $[E_{a|x},F_{b|y}]=0$.
The functional is defined on all finitely many words needed for these
constraints and the observed probabilities; $\ell(W^*)=\overline{\ell(W)}$.
If a moment is absent from the positive matrix it is free apart from the
stated identities. This definition includes standard finite NPA lists;
additional localizing matrices, state- or operator-optimality constraints,
or problem-specific inequalities define other relaxations.

\begin{lemma}
\label{supp:lem:valid-behavior}\label{supp:lem:finite-list-containment}
Every level-$d$ feasible functional with $d\geq2$ defines a valid
correlation through
$E_{a|x}=(e+aa_x)/2$, $F_{b|y}=(e+bb_y)/2$ and
$p(a,b|x,y)=\ell(E_{a|x}F_{b|y})$.
Writing $\mathcal N_d$ for these extracted correlations, every fixed
finite list $\mathcal S$ admits a finite $d\geq2$ such that
$\mathcal N_d\subseteq\mathcal N_{\mathcal S}$.
\end{lemma}
\begin{proof}
$P=E_{a|x}F_{b|y}$ is a self-adjoint projector in $V_2\subseteq V_d$,
so $p=\ell(P)=\ell(P^*P)\geq0$. Summing the projectors gives
$\sum_{ab}p=\ell(e)=1$ and $\sum_bp=\ell(E_{a|x})$, independently
of $y$; Bob's marginal is analogous. This proves validity and preserves
the moments $\langle A_x\rangle_p=\ell(a_x)$,
$\langle B_y\rangle_p=\ell(b_y)$, and
$\langle A_xB_y\rangle_p=\ell(a_xb_y)$.

Expand each projector word using these affine formulas. A word of $D$
projector factors expands into observable words of length at most $D$;
reduction only shortens them. Choose $d\geq2$ at least the largest such
length and enlarge it for any other finitely many required moments. Then
$\bm S=T\bm w_d$ for a real matrix $T$, and
\begin{equation}
\Gamma_{\mathcal S}=T\Gamma_dT^*\succeq0.
\label{supp:eq:congruence}
\end{equation}
The group relations imply all projector identities. Any remaining
unconstrained moments admit a Hermitian linear extension on the reduced-word
basis, assigning conjugate values on inverse pairs. No further positive
matrix involving those moments is required. Hence the extracted correlation
belongs to $\mathcal N_{\mathcal S}$.
\end{proof}

\begin{theorem}
\label{supp:thm:set-separation}\label{supp:cor:local-accumulation}
For every fixed finite projector-word list $\mathcal S$ as above,
$\Q\subsetneq\mathcal N_{\mathcal S}$. Moreover, every neighborhood of
$p_{\rm loc}(a,b|x,y)=\delta_{a,1}\delta_{b,1}$ contains a correlation
in $\mathcal N_{\mathcal S}\setminus\Q$.
\end{theorem}
\begin{proof}
Quantum moments satisfy all the constraints. Choose $d\geq2$ from the
lemma and a primal optimizer at each $\alpha$, extracting
$p_{d,\alpha}\in\mathcal N_d\subseteq\mathcal N_{\mathcal S}$ with
$\beta_\alpha(p_{d,\alpha})=\omega_{\rm NPA}^{(d)}(\alpha)$.
For $\alpha$ sufficiently close to one, Theorem~\ref{supp:thm:fixed-level}
puts this value strictly above $\omega_{\rm Q}(\alpha)$, proving strict
containment. Set
\begin{equation}
e_{d,\alpha}=\omega_{\rm NPA}^{(d)}(\alpha)-\omega_{\rm Q}(\alpha),
\qquad \lambda_\alpha=\frac{2\Delta_\alpha}{\Delta_\alpha+e_{d,\alpha}},
\qquad q_\alpha=(1-\lambda_\alpha)p_{\rm loc}
+\lambda_\alpha p_{d,\alpha}.
\label{supp:eq:mixing-weight}
\end{equation}
Since $e_{d,\alpha}/\Delta_\alpha\to\infty$, the weight is in $(0,1)$
near one and tends to zero. Convexity puts $q_\alpha$ in
$\mathcal N_{\mathcal S}$, whereas
$\beta_\alpha(q_\alpha)=\omega_{\rm L}(\alpha)+2\Delta_\alpha
=\omega_{\rm Q}(\alpha)+\Delta_\alpha>\omega_{\rm Q}(\alpha)$.
Thus $q_\alpha\notin\Q$. Every probability coordinate differs from
$p_{\rm loc}$ by at most $\lambda_\alpha$, proving convergence to the
deterministic correlation.
\end{proof}

\begin{thebibliography}{99}
\small

\bibitem{GigenaEtAl2025}
N.~Gigena, E.~Panwar, G.~Scala, M.~Ara\'ujo, M.~Farkas, and A.~Chaturvedi,
Self-testing tilted strategies for maximal loophole-free nonlocality,
npj Quantum Inf. \textbf{11}, 82 (2025).

\bibitem{BrunnerEtAl2014}
N.~Brunner, D.~Cavalcanti, S.~Pironio, V.~Scarani, and S.~Wehner,
Bell nonlocality, Rev. Mod. Phys. \textbf{86}, 419--478 (2014),
\href{https://doi.org/10.1103/RevModPhys.86.419}{doi:10.1103/RevModPhys.86.419}.

\bibitem{CHSH1969}
J.~F. Clauser, M.~A. Horne, A.~Shimony, and R.~A. Holt,
Proposed experiment to test local hidden-variable theories,
Phys. Rev. Lett. \textbf{23}, 880 (1969).

\bibitem{MiklinEtAl2022}
N.~Miklin, A.~Chaturvedi, M.~Bourennane, M.~Paw\l{}owski, and A.~Cabello,
Exponentially decreasing critical detection efficiency for any Bell inequality,
Phys. Rev. Lett. \textbf{129}, 230403 (2022),
\href{https://doi.org/10.1103/PhysRevLett.129.230403}{doi:10.1103/PhysRevLett.129.230403}.

\bibitem{Eberhard1993}
P.~H.~Eberhard,
Background level and counter efficiencies required for a loophole-free Einstein--Podolsky--Rosen experiment,
Phys. Rev. A \textbf{47}, R747--R750 (1993),
\href{https://doi.org/10.1103/PhysRevA.47.R747}{doi:10.1103/PhysRevA.47.R747}.

\bibitem{Tsirelson1980}
B.~S. Tsirelson,
Quantum generalizations of Bell's inequality,
Lett. Math. Phys. \textbf{4}, 93 (1980).

\bibitem{NPA2007}
M.~Navascu\'es, S.~Pironio, and A.~Ac\'in,
Bounding the set of quantum correlations,
Phys. Rev. Lett. \textbf{98}, 010401 (2007).

\bibitem{AMP2012}
A.~Ac\'in, S.~Massar, and S.~Pironio,
Randomness versus nonlocality and entanglement,
Phys. Rev. Lett. \textbf{108}, 100402 (2012).

\bibitem{BampsPironio2015}
C.~Bamps and S.~Pironio,
Sum-of-squares decompositions for a family of Clauser--Horne--Shimony--Holt-like inequalities and their application to self-testing,
Phys. Rev. A \textbf{91}, 052111 (2015).

\bibitem{SupicBowles2020}
I.~\v{S}upi\'c and J.~Bowles,
Self-testing of quantum systems: a review,
Quantum \textbf{4}, 337 (2020),
\href{https://doi.org/10.22331/q-2020-09-30-337}{doi:10.22331/q-2020-09-30-337}.

\bibitem{Naimark1940}
M.~Neumark,
\href{https://www.mathnet.ru/eng/im3893}{Spectral functions of a symmetric operator},
Izv. Akad. Nauk SSSR Ser. Mat. \textbf{4}, 277--318 (1940).

\bibitem{Fritz2012}
T.~Fritz,
Tsirelson's problem and Kirchberg's conjecture,
Rev. Math. Phys. \textbf{24}, 1250012 (2012),
\href{https://doi.org/10.1142/S0129055X12500122}{doi:10.1142/S0129055X12500122}.

\bibitem{Halmos1969}
P.~R. Halmos,
Two subspaces,
Trans. Am. Math. Soc. \textbf{144}, 381 (1969).

\bibitem{Masanes2005}
Ll.~Masanes,
Extremal quantum correlations for $N$ parties with two dichotomic observables per site,
\href{https://arxiv.org/abs/quant-ph/0512100}{arXiv:quant-ph/0512100} (2005).

\bibitem{PanwarPandyaWiesniak2022}
E.~Panwar, P.~Pandya, and M.~Wie\'sniak,
An elegant scheme of self-testing for multipartite Bell inequalities,
npj Quantum Inf. \textbf{9}, 71 (2023), \href{https://doi.org/10.1038/s41534-023-00735-3}{doi:10.1038/s41534-023-00735-3}.

\bibitem{ScholzWerner2008}
V.~B. Scholz and R.~F. Werner,
Tsirelson's problem,
\href{https://arxiv.org/abs/0812.4305}{arXiv:0812.4305} (2008).

\bibitem{NPA2008}
M.~Navascu\'es, S.~Pironio, and A.~Ac\'in,
A convergent hierarchy of semidefinite programs characterizing the set of quantum correlations,
New J. Phys. \textbf{10}, 073013 (2008).

\bibitem{PironioEtAl2010}
S.~Pironio, M.~Navascu\'es, and A.~Ac\'in,
Convergent relaxations of polynomial optimization problems with noncommuting variables,
SIAM J. Optim. \textbf{20}, 2157--2180 (2010),
\href{https://doi.org/10.1137/090760155}{doi:10.1137/090760155}.

\bibitem{BoydVandenberghe2004}
S.~Boyd and L.~Vandenberghe,
\emph{Convex Optimization} (Cambridge University Press, Cambridge, 2004),
Secs.~5.2.3 and 5.9.1,
\href{https://web.stanford.edu/~boyd/cvxbook/}{author's online edition}.

\bibitem{VandenbergheBoyd1996}
L.~Vandenberghe and S.~Boyd,
Semidefinite programming,
SIAM Rev. \textbf{38}, 49--95 (1996),
\href{https://doi.org/10.1137/1038003}{doi:10.1137/1038003}.

\bibitem{Motzkin1967}
T.~S. Motzkin,
The arithmetic-geometric inequality, in
\emph{Inequalities}, edited by O.~Shisha
(Academic Press, New York, 1967), pp.~205--224.

\bibitem{Hilbert1888}
D.~Hilbert,
\"Uber die Darstellung definiter Formen als Summe von Formenquadraten,
Math. Ann. \textbf{32}, 342 (1888).

\bibitem{Reznick1978}
B.~Reznick,
Extremal PSD forms with few terms,
Duke Math. J. \textbf{45}, 363 (1978).

\bibitem{hs} C. Hanselka and R. Sinn, Positive semidefinite univariate matrix
polynomials, \emph{Math. Z.} \textbf{292}, 83--101 (2019);
\href{https://arxiv.org/pdf/1707.08261}{Corollary 2.8}.
\bibitem{soc} C. Scheiderer, Second-Order Cone Representation for Convex Sets in
the Plane, \emph{SIAM J. Appl. Algebra Geom.} \textbf{5}, 114--139 (2021),
\href{https://doi.org/10.1137/20M133717X}{doi:10.1137/20M133717X};
\href{https://arxiv.org/html/2004.04196v2}{Theorem 0.1}.
\bibitem{prior} A. Chaturvedi, \emph{No Finite NPA Level Characterizes the Complete
Quantum Set in the Simplest Bell Scenario}, project manuscript, September 2026. Its complete argument is reorganized
in Part II of this companion.
\bibitem{Tavakoli2024}
A.~Tavakoli, A.~Pozas-Kerstjens, P.~Brown, and M.~Ara\'ujo, Semidefinite programming relaxations for quantum correlations, Rev. Mod. Phys. \textbf{96}, 045006 (2024), \href{https://doi.org/10.1103/RevModPhys.96.045006}{doi:10.1103/RevModPhys.96.045006}.

\end{thebibliography}
\end{document}
